% GENERATED FILE. Edit canonical lessons, then run npm run generate:books.
% canonical-source-hash: fnv1a64:e4d6d48e1f216860
% canonical-lessons: PA-C103-aktubar, PA-C103-navambar, PA-C103-dasambar, PA-C103-cet, PA-C103-visakh

\chapter{October, November, December, Chet, Visakh}
\label{ch:pa-october-november-december-chet-visakh}
\hlchaptermodality{103}

\begin{chapteropening}
\textbf{By the end of this chapter:} \emph{I can say October, November, December, Chet and Visakh.}

Complete the last lesson of chapter 103: I can say October, November, December, Chet and Visakh.
\end{chapteropening}

\section[\texorpdfstring{\pa{ਅਕਤੂਬਰ} (aktūbar)}{ਅਕਤੂਬਰ (aktūbar)}]{\pa{ਅਕਤੂਬਰ} (aktūbar) — October}

\label{lesson:PA-C103-aktubar}



\begin{quote}
Before the new one: say the Punjabi for August, then the Punjabi for September.
\end{quote}

\subsection*{You'll want to know: \pa{ਅਕਤੂਬਰ}}
\textbf{\pa{ਅਕਤੂਬਰ}} — \emph{aktūbar} — \textquotedblleft{}October\textquotedblright{}.

\begin{cousinweb}
\textbf{\pa{ਅਕਤੂਬਰ}} (\emph{aktūbar}) is a loan from English \emph{October}.
\end{cousinweb}

\begin{grammarlens}[title={What you've built}]
One more word for this chapter.
\end{grammarlens}

\subsection*{Your turn}
\begin{itemize}
\raggedright
  \item \emph{Say it:} \emph{aktūbar}
  \item \emph{Say it:} \emph{aktūbar}, once more
  \item \emph{Say it:} say \emph{aktūbar}
  \item \emph{Recall:} say the Punjabi for August, then the Punjabi for September, then say \emph{aktūbar} again
\end{itemize}

\begin{tcolorbox}[breakable,colback=teal!4,colframe=teal!35!black,title={Before you move on}]
What does \pa{ਅਕਤੂਬਰ} mean? (\textquotedblleft{}October\textquotedblright{}.) Say it once more.
\end{tcolorbox}
\section[\texorpdfstring{\pa{ਨਵੰਬਰ} (navambar)}{ਨਵੰਬਰ (navambar)}]{\pa{ਨਵੰਬਰ} (navambar) — November}

\label{lesson:PA-C103-navambar}



\begin{quote}
Before the new one: say the Punjabi for September, then the Punjabi for October.
\end{quote}

\subsection*{You'll want to know: \pa{ਨਵੰਬਰ}}
\textbf{\pa{ਨਵੰਬਰ}} — \emph{navambar} — \textquotedblleft{}November\textquotedblright{}.

\begin{cousinweb}
\textbf{\pa{ਨਵੰਬਰ}} (\emph{navambar}) is a loan from English \emph{November}.
\end{cousinweb}

\begin{grammarlens}[title={What you've built}]
One more word for this chapter.
\end{grammarlens}

\subsection*{Your turn}
\begin{itemize}
\raggedright
  \item \emph{Say it:} \emph{navambar}
  \item \emph{Say it:} \emph{navambar}, once more
  \item \emph{Say it:} say \emph{navambar}
  \item \emph{Recall:} say the Punjabi for September, then the Punjabi for October, then say \emph{navambar} again
\end{itemize}

\begin{tcolorbox}[breakable,colback=teal!4,colframe=teal!35!black,title={Before you move on}]
What does \pa{ਨਵੰਬਰ} mean? (\textquotedblleft{}November\textquotedblright{}.) Say it once more.
\end{tcolorbox}
\section[\texorpdfstring{\pa{ਦਸੰਬਰ} (dasambar)}{ਦਸੰਬਰ (dasambar)}]{\pa{ਦਸੰਬਰ} (dasambar) — December}

\label{lesson:PA-C103-dasambar}



\begin{quote}
Before the new one: say the Punjabi for October, then the Punjabi for November.
\end{quote}

\subsection*{You'll want to know: \pa{ਦਸੰਬਰ}}
\textbf{\pa{ਦਸੰਬਰ}} — \emph{dasambar} — \textquotedblleft{}December\textquotedblright{}.

\begin{cousinweb}
\textbf{\pa{ਦਸੰਬਰ}} (\emph{dasambar}) is a loan from English \emph{December}.
\end{cousinweb}

\begin{grammarlens}[title={What you've built}]
One more word for this chapter.
\end{grammarlens}

\subsection*{Your turn}
\begin{itemize}
\raggedright
  \item \emph{Say it:} \emph{dasambar}
  \item \emph{Say it:} \emph{dasambar}, once more
  \item \emph{Say it:} say \emph{dasambar}
  \item \emph{Recall:} say the Punjabi for October, then the Punjabi for November, then say \emph{dasambar} again
\end{itemize}

\begin{tcolorbox}[breakable,colback=teal!4,colframe=teal!35!black,title={Before you move on}]
What does \pa{ਦਸੰਬਰ} mean? (\textquotedblleft{}December\textquotedblright{}.) Say it once more.
\end{tcolorbox}
\section[\texorpdfstring{\pa{ਚੇਤ} (cet)}{ਚੇਤ (cet)}]{\pa{ਚੇਤ} (cet) — Chet, the first month of the Nanakshahi calendar}

\label{lesson:PA-C103-cet}



\begin{quote}
Before the new one: say the Punjabi for November, then the Punjabi for December.
\end{quote}

\subsection*{You'll want to know: \pa{ਚੇਤ}}
\textbf{\pa{ਚੇਤ}} — \emph{cet} — \textquotedblleft{}Chet, the first month of the Nanakshahi calendar\textquotedblright{}.

\begin{cousinweb}
\textbf{\pa{ਚੇਤ}} (\emph{cet}) is a tadbhava, a word inherited from Sanskrit and reshaped by sound change. It comes from \emph{caitra}, \textquotedblleft{}the month of the full moon in Citrā\textquotedblright{} (CDIAL 4915). CDIAL lists Punjabi \emph{cetar}, \textquotedblleft{}March-April, spring\textquotedblright{}.
\end{cousinweb}

\begin{grammarlens}[title={What you've built}]
One more word for this chapter.
\end{grammarlens}

\subsection*{Your turn}
\begin{itemize}
\raggedright
  \item \emph{Say it:} \emph{cet}
  \item \emph{Say it:} \emph{cet}, once more
  \item \emph{Say it:} say \emph{cet}
  \item \emph{Recall:} say the Punjabi for November, then the Punjabi for December, then say \emph{cet} again
\end{itemize}

\begin{tcolorbox}[breakable,colback=teal!4,colframe=teal!35!black,title={Before you move on}]
What does \pa{ਚੇਤ} mean? (\textquotedblleft{}Chet, the first month of the Nanakshahi calendar\textquotedblright{}.) Say it once more.
\end{tcolorbox}
\section[\texorpdfstring{\pa{ਵਿਸਾਖ} (visākh)}{ਵਿਸਾਖ (visākh)}]{\pa{ਵਿਸਾਖ} (visākh) — Visakh, the month of Vaisakhi}

\label{lesson:PA-C103-visakh}



\begin{quote}
Before the new one: say the Punjabi for December, then the Punjabi for Chet.

\emph{Read it:} these two lines without stopping — \textbf{\pa{ਮੇਰਾ} \pa{ਨਾਂ} \pa{ਅਮਨ} \pa{ਹੈ।} \pa{ਮੈਂ} \pa{ਪੰਜਾਬੀ} \pa{ਬੋਲਦਾ} \pa{ਹਾਂ।}}
\end{quote}

\subsection*{You'll want to know: \pa{ਵਿਸਾਖ}}
\textbf{\pa{ਵਿਸਾਖ}} — \emph{visākh} — \textquotedblleft{}Visakh, the month of Vaisakhi\textquotedblright{}.

\begin{cousinweb}
\textbf{\pa{ਵਿਸਾਖ}} (\emph{visākh}) is a semi-tatsama: Sanskrit \emph{vaiśākha}, \textquotedblleft{}the month April-May\textquotedblright{}, borrowed after the old sound changes had run their course and then lightly reshaped (CDIAL 12149). CDIAL lists no Punjabi form.
\end{cousinweb}

\begin{grammarlens}[title={What you've built}]
One more word for this chapter.
\end{grammarlens}

\subsection*{Your turn}
\begin{itemize}
\raggedright
  \item \emph{Say it:} \emph{visākh}
  \item \emph{Say it:} \emph{visākh}, once more
  \item \emph{Say it:} say \emph{visākh}
  \item \emph{Recall:} say the Punjabi for December, then the Punjabi for Chet, then say \emph{visākh} again
\end{itemize}

\begin{tcolorbox}[breakable,colback=teal!4,colframe=teal!35!black,title={Before you move on}]
What does \pa{ਵਿਸਾਖ} mean? (\textquotedblleft{}Visakh, the month of Vaisakhi\textquotedblright{}.) Say it once more.
\end{tcolorbox}
